\documentclass[12pt,a4paper]{report}
\usepackage[utf8]{inputenc}
\usepackage[english]{babel}
\usepackage{amsmath}
\usepackage{amsfonts}
\usepackage{amssymb}
\usepackage{graphicx}
\usepackage{hyperref}
\usepackage{geometry}
\usepackage{float}
\usepackage{listings}
\usepackage{xcolor}

\geometry{margin=1in}

\begin{document}

% --- Title Page ---
\begin{titlepage}
    \centering
    {\large \textbf{Università di Modena e Reggio Emilia}} \\
    {\large Dipartimento di Fisica, Matematica e Informatica} \\
    \vspace{2cm}
    {\Large \textit{Computer Graphics Project}} \\
    \vspace{1.5cm}
    \rule{\linewidth}{0.5mm} \\[0.4cm]
    {\huge \textbf{Poisson Image Editing Pro: \\ A Gradient-Domain Manipulation Tool}} \\[0.4cm]
    \rule{\linewidth}{0.5mm} \\
    \vspace{3cm}
    
    \begin{flushright}
        \large
        \textbf{Author:} \\
        [Your Name Here] \\
        \vspace{1cm}
        \textbf{Academic Year 2025-2026}
    \end{flushright}
\end{titlepage}

\tableofcontents
\newpage

% --- Chapter 1: Introduction & Implementation ---
\chapter{Poisson Image Editing}

The project described in this report aims to implement a comprehensive and interactive tool for gradient-domain image editing, based on the seminal framework introduced by Pérez et al. (2003).

\section{Background and Context}
The reference method is that proposed in \textit{"Poisson Image Editing"} (Pérez et al., 2003), which shifts the paradigm of image manipulation from the pixel (intensity) domain to the gradient (variation) domain. 

The core idea is that the human eye is more sensitive to gradients (changes in intensity) than to absolute color values. By solving a Poisson equation with Dirichlet boundary conditions, we can "reconstruct" a source region into a destination image such that the boundaries match perfectly, ensuring a seamless blend regardless of lighting or color differences between the two images.

For more complex scenarios, the project also incorporates:
\begin{itemize}
    \item \textbf{Mixed Gradients:} Combining source and destination variations to preserve background textures (e.g., skin details under a tattoo).
    \item \textbf{Section 4 Applications:} Advanced local modifications including texture flattening, illumination correction, and selective color tinting.
\end{itemize}

\section{Project Implementation}
The theoretical principles were translated into a practical software project written in Python, adhering to the constraint of not using the OpenCV library. The implementation involved the following core activities:

\begin{itemize}
    \item \textbf{Discrete Poisson Solver:} I implemented a vectorized solver using \texttt{NumPy} and \texttt{SciPy}. The tool builds a sparse Laplacian matrix $A$ over the mask region. To ensure scalability, I used a hybrid approach: a direct solver (\texttt{spsolve}) for regions under 100,000 pixels and an iterative Conjugate Gradient (\texttt{cg}) solver for larger selections to prevent memory exhaustion.
    
    \item \textbf{Advanced Editing Modes:}
    \begin{itemize}
        \item \textit{Texture Flattening:} Implemented using a sparse edge sieve via the Sobel operator to keep only high-magnitude gradients.
        \item \textit{Local Illumination:} Implemented dynamic range compression of gradients to reduce specular highlights or bring out shadow details.
        \item \textit{Color Change:} Developed an independent RGB gradient scaling mode to "tint" objects while preserving shading and volume.
    \end{itemize}

    \item \textbf{Professional Interactive UI:} Using \texttt{PySide6}, I developed an all-in-one workspace. Features include:
    \begin{itemize}
        \item Interactive polygon selection.
        \item Real-time positioning and scaling of the selection layer via mouse-drag and wheel-scroll.
        \item Context-aware parameter sliders for different editing modes.
        \item Viewport zoom controls for detailed mask drawing.
    \end{itemize}
    
    \item \textbf{Technical Optimization:} All core math is vectorized to minimize Python loop overhead. Boundary conditions $|N_p|$ are correctly calculated per pixel to handle selections touching image edges accurately.
\end{itemize}

\chapter{Results}

\section{Experimental Evaluation}
The results obtained with the tool are highly satisfactory. Below are placeholders for the experimental results obtained during testing.

% Placeholder for First Result (Cloning)
\begin{figure}[H]
    \centering
    % \includegraphics[width=0.3\linewidth]{source_1}
    % \includegraphics[width=0.3\linewidth]{dest_1}
    % \includegraphics[width=0.3\linewidth]{result_1}
    \vspace{4cm} % Manual space
    \caption{Seamless Cloning: Source region, Destination background, and final Poisson blend.}
\end{figure}

\section{Case Study: Reflection Reduction}
One of the most challenging tasks is the reduction of specular highlights. Using the \textit{Illumination Change} mode, the tool successfully dampens harsh reflections on surfaces like fruit skin by compressing large gradients while preserving the subtle underlying textures.

% Placeholder for Reflection Result
\begin{figure}[H]
    \centering
    % \includegraphics[width=0.45\linewidth]{fruit_original}
    % \includegraphics[width=0.45\linewidth]{fruit_corrected}
    \vspace{5cm} % Manual space
    \caption{Specular Highlight Reduction: Comparison between the original image and the corrected result ($\alpha=0.12, \beta=0.30$).}
\end{figure}

\chapter{Conclusions}
In general, the gradient-domain approach proved extremely effective for both seamless object insertion and local surface editing. The use of the Poisson equation allows for a level of blending that is virtually impossible to achieve with standard pixel-based "feathering" or "alpha-blending."

The implementation of the professional UI significantly lowers the barrier to entry for using these mathematical techniques, providing a workflow that feels natural and responsive. Future improvements could involve GPU-accelerated solvers (e.g., via \texttt{PyTorch} or \texttt{CuPy}) to handle ultra-high-resolution images in real-time, although the current iterative CPU solver provides a robust and portable baseline.

\end{document}
